% ============================================================
%  Closing the Loop from Measured Activity to Cell-Resolved Response:
%  A Computational Framework for Spatially Explicit PRRT Modeling
%  Hunter Kinder & Brett Faulkner
%  Preprint - not peer reviewed - September 2026
%  This manuscript reports no live culture, recombinant/synthetic
%  nucleic-acid experiment, human or animal research, or physical
%  irradiation experiment. It is a computational framework
%  specification and verified literature-gap analysis; it reports no
%  executed simulation results.
% ============================================================
\documentclass[11pt]{article}

% --- Geometry -----------------------------------------------
\usepackage[top=1in,bottom=1in,left=1.1in,right=1.1in]{geometry}

% --- Encoding & fonts ----------------------------------------
\usepackage[T1]{fontenc}
\usepackage[utf8]{inputenc}
\usepackage{lmodern}
\usepackage{microtype}

% --- Math ---------------------------------------------------
\usepackage{amsmath,amssymb,amsthm}
\usepackage{bm}

% --- Tables -------------------------------------------------
\usepackage{booktabs}
\usepackage{array}
\usepackage{tabularx}
\usepackage{longtable}

% --- Graphics & colour ----------------------------------------
\usepackage{graphicx}
\usepackage{xcolor}
\definecolor{linkblue}{RGB}{0,70,150}

% --- Hyperref ------------------------------------------------
\usepackage[colorlinks=true,
            linkcolor=linkblue,
            citecolor=linkblue,
            urlcolor=linkblue,
            pdftitle={Closing the Loop from Measured Activity to Cell-Resolved Response: A Computational Framework for Spatially Explicit PRRT Modeling},
            pdfauthor={Hunter Kinder, Brett Faulkner}]{hyperref}

% --- Layout ---------------------------------------------------
\usepackage{parskip}
\setlength{\parskip}{6pt}
\usepackage{setspace}
\setstretch{1.15}
\usepackage{titlesec}
\titleformat{\section}{\large\bfseries}{\thesection.}{0.5em}{}
\titleformat{\subsection}{\normalsize\bfseries}{\thesubsection}{0.5em}{}
\titlespacing*{\section}{0pt}{14pt}{6pt}
\titlespacing*{\subsection}{0pt}{10pt}{4pt}

% --- Floats & captions ----------------------------------------
\usepackage{caption}
\captionsetup{font=small,labelfont=bf}

% --- Line breaking --------------------------------------------
\emergencystretch=3em

% ============================================================
\begin{document}

% ---- Title block --------------------------------------------
\begin{center}
  {\LARGE\bfseries Closing the Loop from Measured Activity to\\
   Cell-Resolved Response}\\[10pt]
  {\large\bfseries A Computational Framework Specification and Verified Gap\\
   Analysis for Spatially Explicit Modeling of Peptide Receptor\\
   Radionuclide Therapy in Neuroendocrine Tumors}\\[14pt]
  {\normalsize
    \textbf{Hunter Kinder}, B.A., M.A.T.L. --- Independent Researcher, Missouri, USA\\[2pt]
    \textbf{Brett Faulkner}, B.Sc. --- Independent Researcher\\[10pt]
  }
  {\small
    \textit{Preprint --- not peer reviewed}\\
    \textit{Version 1.0 --- September 2026}
  }
\end{center}
\vspace{4pt}
\noindent\rule{\linewidth}{0.5pt}
\vspace{2pt}

% ---- Abstract ------------------------------------------------
\begin{abstract}
\noindent
Peptide receptor radionuclide therapy (PRRT) fails in the clinic through a spatial process:
somatostatin receptor subtype 2 (SSTR2) expression varies within and between lesions, the
delivered absorbed dose is heterogeneous at the cellular scale, and retreatment appears to
select for low-expression clones. Each fragment of this process is measured and modelled
somewhere in the literature. No study joins the fragments: small-scale dosimetry stops at dose
histograms, biological response models assume uniform dose or abstract geometry, and no
published work couples a measured intratumoral activity distribution to a spatially explicit,
cell-resolved response model that feeds selection dynamics back into the next treatment cycle.
We verified that absence by systematic search against explicit refutation criteria, and we catalog the near-miss studies
that fail individual criteria. The searched sources were EuropePMC, PubMed Central, OpenAlex, Crossref, CORE, arXiv, and
alphaXiv (September 2026).

This manuscript specifies the missing model and audits what the specification does and does not
support. The framework initializes a Cellular Potts Model (CPM) from same-section quantitative autoradiography and
histology. It transports electron and alpha dose through the evolving cell geometry, drives cell-state transitions with
protracted-exposure linear-quadratic kinetics, and carries a heritable SSTR2 expression trait. That trait's differential
survival under heterogeneous dose closes the loop into the next cycle's activity map. The one-way dose-transfer foundation
(geometry export, transport, dose conversion, cell-level attribution) is inherited from our
validated CPM--OpenMC coupling and is reused here unchanged; what is new is the response layer
and the activity-feedback specification. No simulation results are reported: every executed
quantity in this preprint is either a literature prior or a verified absence.

The boundaries are declared rather than discovered late. The parameters are separated by
provenance (measured inputs, literature priors, declared modelling choices) and are not reported
as though they were uniform in evidentiary weight. The model can test clone selection under
heterogeneous dose; it cannot represent immune response, bystander signalling, or vascular
remodelling, and it makes no patient-specific prediction until parameterised against
patient data. What the work establishes is a reproducible, fully specified loop and a verified
statement of the gap it fills, with the measurements and controls required before any clinical
claim can be evaluated.
\end{abstract}
\noindent\textbf{Keywords:} peptide receptor radionuclide therapy; neuroendocrine tumors;
SSTR2 heterogeneity; cellular Potts model; microdosimetry; tumor control probability;
digital twin
\par
\noindent\rule{\linewidth}{0.5pt}

% ---- 1. Introduction -----------------------------------------
\section{Introduction}\label{sec:intro}

Peptide receptor radionuclide therapy delivers radiation to neuroendocrine tumor (NET) and
pheochromocytoma/paraganglioma (PPGL) lesions through somatostatin-receptor targeting of
beta-emitting lutetium-177 or alpha-emitting actinium-225 conjugates; in advanced PPGLs
pooled disease control reaches 90\% with a pooled median progression-free survival near
32 months \cite{su2023ppgl}. The therapy works, and
its limits are known. Lesions that respond to a first cycle can progress under later ones,
and the clinical pathology literature locates part of the failure inside the tumor's own
spatial structure. SSTR2 expression differs between patients, between lesions, and between
clones within one lesion. DNA damage after therapy concentrates in high-expression regions
while low-expression regions are spared, and tumors that recur after treatment carry less
SSTR2 than the treatment-naive primary \cite{feijtel2021,fransson2024}. A pediatric
neuroblastoma case treated with $^{177}$Lu-DOTATATE shows the same pattern at clone
resolution: a metastasis with weak SSTR2 staining absorbed 2~Gy where the better-expressing
primary absorbed 54~Gy, and the lesion that progressed was the weakly expressing one
\cite{fransson2024}.

The dosimetry literature measures the other half of the story. Digital autoradiography of
cryosections, registered to hematoxylin-and-eosin-stained neighbours, resolves the
intratumoral activity distribution at tens of micrometres, and Monte Carlo transport through
that measured source gives absorbed-dose distributions that a mean-organ or even a
voxol-level model cannot recover \cite{mellhammar2023}. The same workflow extends to
alpha emitters, where autoradiography and histology are acquired on the same section and
stacked into a three-dimensional activity volume for microscale dosimetry \cite{sahota2025}.
Cellular-scale dosimetry of $^{177}$Lu-DOTATATE shows that cluster geometry and subcellular
source location move nuclear absorbed dose by factors of two or more
\cite{tamborino2020,tamborino2025}, and heterogeneous targeting leaves unlabelled cells at
3--17\% of the dose they would receive under uniform targeting \cite{larouze2023}.

Between these two literatures sits a gap that is easy to state and hard to excuse. Dosimetry
papers stop at the dose histogram and a tumor control probability (TCP) computed from it.
Biological response models start from a uniform dose or an abstract geometry and evolve
cell populations forward. Neither feeds the measured heterogeneity into a cell-resolved
spatial model and asks what the selection dynamics do, cycle after cycle, to the expression
distribution that the next dose will see. Section~\ref{sec:gap} verifies that no published
study closes that loop, and catalogues the studies that come closest.

This manuscript does three things. First, it states the closed-loop framework completely enough to be executed
(Section~\ref{sec:framework}): histology-initialized CPM geometry, cell-resolved dosimetry,
protracted-exposure response kinetics, and a heritable SSTR2 trait whose differential
survival feeds the next cycle's activity map. Second, it separates every parameter by provenance ---
measured input, literature prior, or declared modelling choice --- and refuses to report
them as though they carried the same evidentiary weight (Section~\ref{sec:params}).
Third, it declares the model's boundaries before any reader has to ask: what the
specification can test, what it cannot represent, and what must be measured before a
clinical claim is even formulable (Section~\ref{sec:boundaries}).

The position of this work relative to our prior preprint is simple. There we coupled a CPM
to OpenMC photon transport, executed the one-way path end to end, and audited what the
coupling did and did not support \cite{kinder2026biofilm}. The dose-transfer machinery ---
geometry export, transport, dose conversion, cell-level attribution with orientation and
conservation checks --- is reused unchanged. What was missing there, and is specified here,
is the biological response layer and the feedback that makes the activity map a function of
what the previous cycle did to the cell population.

% ---- 2. Related Work -----------------------------------------
\section{Related Work: The Broken Chain}\label{sec:gap}

\subsection{Measured-activity microdosimetry}

Mellhammar et al.\ computed cell-weighted absorbed-dose distributions for $^{177}$Lu,
$^{90}$Y, and $^{225}$Ac by defining the Monte Carlo source from digital autoradiography of
xenograft cryosections and the cell positions from hematoxylin-eosin segmentation
\cite{mellhammar2023}. Their result is the standing demonstration that voxel dose and cell
dose are different quantities: low-dose voxels that contain few cells contribute less to the
cell-dose histogram than their volume suggests, and TCP estimates built on the two
distributions differ materially. Sahota et al.\ integrated quantitative autoradiography and
histological imaging on the same tissue section for alpha-emitter therapy, using the
anatomical image to stack sections into a three-dimensional activity volume for microscale
dosimetry \cite{sahota2025}. Tamborino et al.\ built image-based multicellular dosimetry for
$^{177}$Lu-DOTATATE in cluster-forming cell lines, showing cross-dose sensitivity to real
cell placement that idealized packing models underestimate by 66--90\%
\cite{tamborino2025}, and earlier established the sensitivity of nuclear S-values to
measured cell morphology and to subcellular source location, including a Golgi-localized
source raising nuclear dose by up to 149\% relative to uniform cytoplasmic activity
\cite{tamborino2020}. Larouze et al.\ quantified what heterogeneous targeting costs in a
19-cell cluster: with four cells unlabelled, membranes of unlabelled cells receive on
average 9.6\% ($^{177}$Lu) or 2.9\% ($^{161}$Tb) of the uniform-targeting dose
\cite{larouze2023}. None of these studies evolves a cell population; the dosimetry output
is a distribution, and the biology is a static density map.

\subsection{Biological response models without measured input}

Liu, Swat, Glazier et al.\ built an agent-based cell-state transition model for
post-irradiation response, validated against monolayer irradiation outcomes and framed as
a digital-twin component \cite{liu2025pbt}. It represents cell phase evolution, phenotype
evolution, and survival after irradiation, but its dose input is external-beam and
effectively uniform; it carries no radionuclide, no receptor heterogeneity, and no measured
activity distribution. Cordoni et al.\ extended the generalized stochastic microdosimetric
model (GSM$^2$) to an agent-based framework with cell-cycle dynamics, dose-rate dependence,
and fractionation effects on a fixed lattice \cite{cordoni2026}; again the radiation field
is prescribed, not measured, and the application is particle therapy rather than
radiopharmaceuticals. Albanese et al.\ modelled spatially resolved oxygenation and
continuous phenotypic adaptation in a trait-structured partial-differential-equation system
and used it to rank fractionation schedules \cite{albanese2026}; the framework is
cell-resolved in neither geometry nor dose. Mvogo et al.\ compared $^{90}$Y, $^{177}$Lu,
$^{131}$I, and $^{225}$Ac extrapolated dose rates in a delay-differential cancer-stem-cell
model with linear-quadratic kinetics \cite{mvogo2026}; the model has no space at all. Zhang
et al.\ resolved the peroxyl-radical and NRF2/GSH kinetics behind FLASH sparing in a
two-compartment well-mixed ODE system \cite{zhang2025flash} --- a chemical-field kinetics
template this framework's ROS layer would inherit, once such a layer exists.

\subsection{The transport bridge}

RADCELL couples Geant4 (with Geant4-DNA low-energy physics) to CompuCell3D: tumor cellular
geometry from the CPM is exported to Geant4, dose and DNA-damage distributions are
transported through it, and the results are fed back per cell \cite{liu2021radcell}. This
is the existence proof that a CPM$\leftrightarrow$transport bridge can be built, and our
own OpenMC--CPM coupling validated the same bridge class for photons with conservation and
orientation checks across the exchange boundary \cite{kinder2026biofilm}. Both bridges
were exercised on synthetic geometries; neither has been driven by a measured
histology-initialized activity map.

\subsection{The verified gap}

Ryhiner et al.\ set out the digital-twin roadmap for radiopharmaceutical therapy and name
patient-specific biological response as a required, still-missing component
\cite{ryhiner2025tdt}. The claim of this manuscript is narrower and was tested rather than
asserted:

\begin{quote}
\textbf{Gap claim.} No published study couples a \emph{measured} intratumoral activity
distribution to a \emph{spatially explicit, cell-resolved} response model for PRRT in which
dose-driven cell dynamics feed back into the subsequent activity distribution.
\end{quote}

A candidate would have to satisfy three criteria simultaneously. (a) A measured or realistic heterogeneous activity
input, not an assumed-uniform source. (b) A cell-resolved spatial model --- CPM, individual-based model, or cellular
automaton --- not a continuum PDE/ODE system and not a dosimetry-only pipeline. (c) A closed response loop, where dose
changes the cell population and the changed population changes the next dose calculation. We
searched EuropePMC, PubMed Central, OpenAlex, Crossref, CORE, bioRxiv, medRxiv, arXiv, and
alphaXiv with PRRT- and model-class-keyed query vectors (September 2026), triaged all
returned records, and characterized each near-miss in full text. PubMed returned a parse
error on the unified query and arXiv rate-limited two queries; EuropePMC, which indexes
PubMed full text, carried that coverage, and alphaXiv discovery covered arXiv content. No
candidate satisfies all three criteria. Table~\ref{tab:nearmiss} catalogues the closest
studies and the criterion each fails.

\begin{table}[t]
\centering
\caption{Near-miss studies and the refutation criterion each fails. Criteria: (a)
measured heterogeneous activity input; (b) cell-resolved spatial model; (c) closed
response loop into the next activity distribution.}
\label{tab:nearmiss}
\small
\begin{tabularx}{\linewidth}{@{}l X c c c@{}}
\toprule
Study & What it does & (a) & (b) & (c) \\
\midrule
\cite{mellhammar2023} & Measured-source MC dosimetry; cell-weighted dose and TCP &
yes & no & no \\
\cite{sahota2025} & Same-section autoradiography--histology 3D activity volume &
yes & no & no \\
\cite{tamborino2025} & Image-based multicellular dosimetry in clusters &
no & partly & no \\
\cite{larouze2023} & Heterogeneous targeting in a fixed 19-cell cluster &
no & partly & no \\
\cite{liu2025pbt} & Agent-based post-irradiation cell-state transitions &
no & yes & no \\
\cite{cordoni2026} & Agent-based GSM$^2$ extension, fractionation on a lattice &
no & yes & no \\
\cite{liu2021radcell} & Geant4--CC3D bridge, synthetic geometry &
no & yes & no \\
\cite{kinder2026biofilm} & CPM--OpenMC one-way coupling, validated, synthetic geometry &
no & yes & no \\
\cite{albanese2026} & Oxygenation and phenotype adaptation, trait-structured PDE &
no & no & no \\
\cite{mvogo2026} & Radionuclide dose-rate comparison, ODE/DDE stem-cell model &
no & no & no \\
\bottomrule
\end{tabularx}
\end{table}

The pieces exist. The loop has never been closed.

% ---- 3. Computational Framework ------------------------------
\section{Computational Framework}\label{sec:framework}

\subsection{Scope and execution state}

The framework is a specification. Nothing in Section~\ref{sec:framework} has been executed
as a coupled simulation in this preprint, and no result of this manuscript depends on the
loop running. What has been executed and validated, in our prior work, is the one-way
transport foundation: geometry export from a CPM, Monte Carlo transport, dose conversion,
and cell-level attribution, with conservation and orientation checks across the exchange
boundary \cite{kinder2026biofilm}. The response layer and activity feedback specified here
are the additions that make the loop closable, and Section~\ref{sec:params} states which
of their parameters are measured, which are literature priors, and which are declared
modelling choices.

\subsection{Notation}

Positions $\bm{x} \in \Omega \subset \mathbb{R}^3$ label voxels of the reconstructed
tissue volume (the radiohistology product). Cells are indexed $i$, occupy voxel sets
$C_i \subset \Omega$, and carry state $(\phi_i, e_i, V_i, S_i)$: cell-cycle phase,
SSTR2 expression level, volume, and surface. Time $t$ is biological time; treatment
cycles are indexed $k$. Activity per voxel at the start of cycle $k$ is $A^{(k)}_j$
[Bq], with time-integrated activity $\tilde{A}^{(k)}_j = \int A^{(k)}_j(t)\,dt$ [Bq\,s].

\subsection{Radiohistology initialization}\label{sec:init}

The model is initialized from three registered, same-section or same-volume maps:

\begin{enumerate}
\item \textbf{Activity} $A_j$, from quantitative same-section autoradiography with
histological context, decay-corrected and normalized per the workflow of Sahota et al.\
\cite{sahota2025}, or from serial-section autoradiography registered to histology per
Mellhammar et al.\ \cite{mellhammar2023}.
\item \textbf{Cell density} $n_j$ and nuclear positions, from hematoxylin-eosin
segmentation, giving viable-cell, necrotic, and stromal masks
\cite{mellhammar2023,sahota2025}.
\item \textbf{Expression} $e_i$, from SSTR2 immunohistochemistry, mapped to the cell
lattice; where IHC is unavailable, the initialization draws $e_i$ from a lognormal
distribution whose parameters are declared priors (Section~\ref{sec:params}).
\end{enumerate}

Cells are placed by sampling the nuclear-position map; each cell's initial CPM volume is
grown to its target volume under the Potts dynamics before cycle 1 begins. The
registration step is not specified here beyond its inputs and its tolerance: elastic
registration of histology to the tissue volume follows established practice
\cite{albers2021}, and a misregistration tolerance enters the uncertainty budget rather
than the dynamics.

\subsection{Cell-resolved dosimetry through the evolving geometry}

For cycle $k$, the absorbed dose to cell $i$ is split into self- and cross-terms:
\begin{equation}
\dot{D}^{(k)}_i = \underbrace{S_{\text{self}}(e_i)\,\tilde{a}^{(k)}_i}_{\text{self}}
\;+\; \sum_{j \neq \text{cells}(i)} S(\bm{r}_i \leftarrow \bm{r}_j)\,
\tilde{A}^{(k)}_j ,
\label{eq:doserate}
\end{equation}
where $\tilde{a}^{(k)}_i$ is the time-integrated activity in cell $i$ and $S$-values are
computed either by Monte Carlo electron transport through the exported geometry ---
which requires an electron-transport code of the Geant4/GATE class, validated for
$^{177}$Lu beta spectra against EGSnrc to 1.2--1.5\% \cite{perrot2014} --- or by
pre-computed cellular S-values for the self term \cite{tamborino2020}, with $^{177}$Lu
betas of 497~keV endpoint and mean penetration near 0.2~mm in tissue
\cite{ijms2026nanoparticle}. The OpenMC path inherited from our prior work transports
photons and neutrons only \cite{kinder2026biofilm}; it cannot serve as the reference
arm for the beta dose, and the convergence study of Section~\ref{sec:params} must
name its reference code accordingly. For $^{225}$Ac chains the
alpha range of 50--100~$\mu$m makes the self- and near-neighbour terms dominate
\cite{sahota2025}. The transport is executed on the \emph{current} cycle geometry; the
activity map $A^{(k)}$ is attached to voxels, not cells, so cell motion and death between
transport calls are handled by re-attribution at the next call.

Equation~\ref{eq:doserate} is a specification, not a method: the numerical choice between
Monte Carlo re-transport per cycle and cached dose-point-kernel lookups is a convergence
question that Section~\ref{sec:params} assigns to the uncertainty budget, and the
specification refuses to fix it without that study.

\subsection{Response: protracted-exposure survival kinetics}

PRRT delivers dose at a low rate over days, so repair during irradiation cannot be
ignored. Cell $i$'s survival to the end of cycle $k$ is
\begin{equation}
\mathrm{SF}^{(k)}_i = \exp\!\left( -\alpha_{\phi}\, D^{(k)}_i
\;-\; \beta_{\phi}\, G(T)\, \bigl(D^{(k)}_i\bigr)^2 \right),
\label{eq:sf}
\end{equation}
with the Lea--Catcheside factor for a constant dose rate over active duration $T$ and
monoexponential repair of half-time $T_{\text{rep}}$,
\begin{equation}
G(T) = \frac{2}{(\mu T)^2}\left( \mu T - 1 + e^{-\mu T} \right),
\qquad \mu = \ln 2 / T_{\text{rep}} .
\label{eq:lea}
\end{equation}
Radiosensitivity $\alpha_{\phi}$ and $\beta_{\phi}$ carry a cell-cycle-phase dependence
(G2/M most sensitive), and $\alpha_{\phi}$ is further modulated by oxygenation where an
oxygen map is available; both dependences are literature priors with declared ranges
(Section~\ref{sec:params}). Death is realized in the CPM as a stochastic transition at
mitosis entry: a cell that has accumulated cycle-$k$ dose attempts division with
probability $\mathrm{SF}^{(k)}_i$ and otherwise enters a dying state, is cleared by
volume resorption, and leaves the lattice. This places lethality at mitotic catastrophe,
consistent with the proliferative-dependence of radiation death in spheroid and solid
tumor models, and it couples response to the growth dynamics the CPM already carries.

\subsection{The heritable expression trait and the closed loop}

SSTR2 expression $e_i$ is a continuous heritable trait. Uptake attaches activity to cells
in proportion to expression:
\begin{equation}
A^{(k+1)}_j \;\propto\; \rho^{(k+1)}_j \cdot \bar{e}^{(k+1)}_j ,
\label{eq:feedback}
\end{equation}
where $\rho^{(k+1)}_j$ is the viable-cell density and $\bar{e}^{(k+1)}_j$ the mean
expression of the cells occupying voxel $j$ at the start of cycle $k{+}1$, normalized to
the administered activity and the measured whole-lesion time-integrated activity.
Equation~\ref{eq:feedback} is the loop-closing term and is a \emph{declared modelling
choice}, not a measurement: it encodes the first-order approximation that uptake follows
receptor density, and it is the point of the model. Because high-$e$ cells receive more
cross-dose (Eq.~\ref{eq:doserate}), die more (Eq.~\ref{eq:sf}), and are therefore
under-represented in $\bar{e}^{(k+1)}$, the population's expression distribution drifts
down under successive cycles --- the retreatment-selection mechanism observed clinically
\cite{feijtel2021,fransson2024} --- and the next cycle's dose follows it. Daughter cells
inherit $e$ with a small perturbation of declared variance; the perturbation is a
standing variation parameter, and Section~\ref{sec:params} fixes its scan range.

With Eq.~\ref{eq:feedback} removed, the model reduces to the open-loop configuration
(measured $A^{(k)}$ replayed for every cycle), which is the natural control experiment:
the divergence between the open-loop and closed-loop trajectories \emph{is} the
retreatment-selection prediction.

\subsection{Potts dynamics and the Hamiltonian}

The tissue mechanics are standard Glazier--Graner--Hogeweg dynamics with Hamiltonian
\begin{equation}
H = \sum_i \lambda_V \bigl(V_i - V_i^{\text{target}}\bigr)^2
  + \sum_i \lambda_S \bigl(S_i - S_i^{\text{target}}\bigr)^2
  + \sum_{\text{nn}} J\bigl(\tau_i, \tau'_j\bigr)\bigl(1 - \delta_{\sigma_i \sigma_j}\bigr)
  + H_{\text{growth}},
\label{eq:hamiltonian}
\end{equation}
over nearest-neighbour copy attempts with Metropolis acceptance; cell types $\tau$
include viable tumor cells by clone, dying cells, necrotic debris, and stroma.
$H_{\text{growth}}$ drives target volume toward doubling over the clone-specific
doubling time, at which point division is attempted and Eq.~\ref{eq:sf} is applied.
No chemotaxis, adhesion-modification, or force-coupling to dose is included; the
specification adds no mechanism beyond what the response layer requires, and every
addition would have to be declared with its provenance as these are.

\subsection{What the loop measures}

The experimental outputs of the specification, once executed, are: (i) the
divergence of TCP between the uniform-dose assumption and the histology-initialized
heterogeneous reality, per lesion; (ii) the trajectory of the expression distribution
$\{e_i\}$ over simulated cycles, and the conditions under which low-expression clones
come to dominate; (iii) the sensitivity of (i) and (ii) to the declared choices
(Eq.~\ref{eq:feedback} proportionality, uptake-lognormal spread, death timing). Output
(i) is the model-side answer to the voxel-dose-versus-cell-dose gap demonstrated by
Mellhammar et al.\ \cite{mellhammar2023}; output (ii) is the model-side answer to the
retreatment-failure observations \cite{feijtel2021,fransson2024}; output (iii) is what
makes the other two publishable rather than anecdotal.

\subsection{Data export architecture}

Every executed study writes the spatial state at each treatment cycle as a
regular voxel volume: cell identifier, tissue class, local activity,
absorbed dose, survival fraction, and expression trait, one array per
quantity, with the lattice spacing declared in the file metadata. The
volumes are exchanged as a wrapper index over the per-cycle files, so the
time dimension is navigable in standard scientific-visualization tooling.
The format is a deliberate exchange architecture, not a visualization
convenience: continuum field couplings attach to the same grid without
re-layout. The level-set treatment of cell membranes
\cite{shen2022diffuse} and the pressure-carrying tumor-growth system of
\cite{li2026chns} are the declared targets such a coupling would follow.
What the exported series contains is a time-dependent histology-dose state:
no velocity or pressure field is computed, and none is exported.

% ---- 4. Parameter provenance ---------------------------------
\section{Parameter Provenance}\label{sec:params}

The parameters are not uniform in provenance and are not reported as though they were.
Table~\ref{tab:provenance} separates them. Three rules govern the table: a measured
input is only a quantity read off the radiohistology product for the simulated lesion; a
literature prior carries a citation and an adopted range but was not measured for this
system; a declared modelling choice is a specification decision this manuscript makes
without external evidence and must be scanned, not fixed. No parameter is fitted to any
outcome of the model itself; the framework is explicitly non-calibrating in the sense of
our prior preprint \cite{kinder2026biofilm}.

\begin{table}[t]
\centering
\caption{Parameter provenance. M = measured input; L = literature prior; D = declared
modelling choice. Ranges for L and D entries are adoption or scan ranges, not fits.}
\label{tab:provenance}
\small
\begin{tabularx}{\linewidth}{@{}l l X l@{}}
\toprule
Quantity & Class & Source / justification & Ref \\
\midrule
Activity map $A^{(1)}_j$ & M & Same-section autoradiography, decay-corrected &
\cite{sahota2025,mellhammar2023} \\
Cell density / masks $n_j$ & M & H+E segmentation of the same volume &
\cite{mellhammar2023} \\
SSTR2 map (per-cell $e_i$) & M$|$D & IHC where available; lognormal draw otherwise &
\cite{feijtel2021} \\
$S$-values (self term) & L & Cellular dosimetry, measured morphologies &
\cite{tamborino2020} \\
$^{177}$Lu spectrum / ranges & L & Decay data; mean $\beta$ range $\approx 0.2$~mm &
\cite{ijms2026nanoparticle} \\
$^{225}$Ac alpha range & L & 50--100~$\mu$m chain alphas & \cite{sahota2025} \\
$\alpha_{\phi}, \beta_{\phi}$ (NET cells) & L & Protracted LQ for neuroendocrine lines &
\cite{tamborino2025} \\
Repair half-time $T_{\text{rep}}$ & L & LQ fits under PRRT exposure &
\cite{tamborino2025} \\
Oxygen effect on $\alpha$ & L & Where an oxygen map exists; else omitted &
--- \\
Uptake $\propto e_i$ (Eq.~\ref{eq:feedback}) & D & First-order binding approximation &
--- \\
Uptake lognormal spread $\sigma_e$ & D & Scan range declared &
\cite{feijtel2021,fransson2024} \\
Expression drift at division & D & Standing variation; scan range declared &
--- \\
Misregistration tolerance & D & Enters uncertainty budget, not dynamics &
\cite{albers2021} \\
Transport caching vs.\ re-transport & D & Study complete: DPK inadmissible at
boundaries (cell p95 = 23.8\%); full re-transport required & \cite{kinder2026biofilm} \\
CPM $\lambda_V, \lambda_S, J$ & L$|$D & Priors from CPM literature; final values
declared & --- \\
\bottomrule
\end{tabularx}
\end{table}

Two entries deserve their own sentence. The uptake proportionality
(Eq.~\ref{eq:feedback}) is the load-bearing declared choice: it is what makes the loop
close, it is not directly measured for a cell population in vivo, and every conclusion of
the executed model must be stated conditional on it. The SSTR2 map is the one quantity
that can be measured per-lesion today; where it is measured, the model's initialization is
fully empirical, and where it is drawn, the initialization is a declared prior and the
sensitivity study is mandatory.

% ---- 5. Declared boundaries ----------------------------------
\section{Declared Boundaries}\label{sec:boundaries}

The model's phenotype boundary is selection under heterogeneous dose. It can test whether
measured heterogeneity, transported realistically and responded to with protracted LQ
kinetics, drives the expression distribution toward the low-SSTR2 clones that
retreatment failure implicates \cite{feijtel2021,fransson2024}. It cannot test:

\begin{itemize}
\item \textbf{Immune response.} No immune cell type is in the state space. PRRT's immunological interactions are
outside the model.
\item \textbf{Bystander signalling.} All dose is transported; no unirradiated cell
receives a signal from an irradiated neighbour.
\item \textbf{Vascular remodelling and pharmacokinetics.} The activity map is an input
per cycle; the model does not compute it from vasculature, and inter-cycle changes enter
only through Eq.~\ref{eq:feedback}.
\item \textbf{Interstitial flow and pressure.} No fluid field is specified
and none is exported. Interstitial pressure enters the tumor-growth
literature as an explicit continuum field \cite{li2026chns}. The natural
extension treats cell boundaries as phase-field level sets coupled to an
incompressible Navier--Stokes solver \cite{shen2022diffuse}. The exported
per-cycle voxel arrays are the substrate such a coupling would consume.
\item \textbf{Electron transport.} The inherited one-way transport path (OpenMC) moves
photons and neutrons, not primary electrons \cite{kinder2026biofilm}. Every beta-dose
reference calculation therefore requires a Geant4/GATE-class electron code; a dose-point
kernel cached from such a code is the approximation under test, never the ground truth.
\item \textbf{Sub-cellular dose.} Dose is attributed to cells, not nuclei; the
S-value literature shows nuclear dose depends on subcellular source location
\cite{tamborino2020}, so the model inherits that as declared uncertainty, not resolved
physics.
\item \textbf{ROS chemistry.} No reactive-oxygen-species field is specified. The FLASH
kinetics literature \cite{zhang2025flash} is the template a future chemistry layer would
adapt; until it exists, the model speaks only through LQ response.
\item \textbf{Patient-specific prediction.} Until every D-class parameter is either
measured or bounded by a sensitivity study, no output is a prediction about a patient.
The model predicts relative divergence between assumptions, not absolute response.
\end{itemize}

Two thresholds are permanently unset, following the rule that a prospective decision
parameter must never be chosen from the data that would certify it
\cite{kinder2026biofilm}: no clinical action threshold on TCP divergence, and no
expression level below which a clone is declared resistant. The model may compute such
quantities; it may not recommend them.

% ---- 6. Conclusion -------------------------------------------
\section{Conclusion}

This preprint specifies a closed-loop, histology-initialized, cell-resolved model of PRRT response in neuroendocrine
tumors. It verifies by systematic search that no published work couples measured activity distributions to spatially
explicit selection dynamics. It separates every parameter by provenance so the specification cannot be mistaken for a
calibration. It reports no simulation results. The one-way transport foundation is
inherited from validated prior work \cite{kinder2026biofilm}; the response layer and the
activity-feedback term are specified here for the first time, and the near-miss catalogue
(Table~\ref{tab:nearmiss}) states exactly which existing studies each piece builds on.

What is required before the loop's first output can be believed: a measured
activity-plus-histology-plus-IHC volume for at least one lesion (inputs that exist today
\cite{sahota2025,mellhammar2023}), the transport-caching convergence study, the declared
sensitivity scans, and a demonstration that the open-loop control reproduces the
voxel-versus-cell TCP divergence already established for static dosimetry
\cite{mellhammar2023}. What is required before any clinical claim: those, plus an
independent measurement of the uptake--expression relation that Eq.~\ref{eq:feedback}
currently declares. The framework is offered as a reproducible specification for exactly
that programme.

Three studies in that programme are now executed, each against a pre-declared
protocol with negative-control gates, and each committed with its protocol,
code, and verdict.

Study~1 (transport caching) tested the cheap shortcut. The cached
dose-point kernel was applied to the hashed synthetic geometry by
convolution, with the kernel's dose spread evaluated by fast-Fourier-
transform (FFT) convolution. Against a 128M-decay Geant4 re-transport
reference it fails the pre-declared cell-level criterion of 2\%:
cell p95 $= 23.8\%$. The failure is an escape-surface shell one to two
voxels thick, caused by the convolution's infinite-medium assumption.
Every treatment cycle therefore re-transports.

Study~2 (open-loop response) verified the response layer. The protracted
linear-quadratic kinetics of Eq.~\ref{eq:sf} with the Lea--Catcheside
factor of Eq.~\ref{eq:lea}, applied to the validated dose field, reproduces
the Mellhammar tumor-control divergence at equal total dose. The mechanism
is the inequality of Jensen, the convexity bound on a mean, acting on the
convex region of the protracted survival function. The divergence
exceeds the transport-noise floor by three orders of magnitude.

Study~3 (closed loop) closed the loop. Four simulated treatment cycles
applied death at mitosis and re-uptake proportional to the heritable
expression trait. The population's mean log-expression falls by $0.48$.
An open-loop control that freezes the activity map falls only $0.19$.
The selection contrast is the model-side reproduction of the
retreatment-failure pattern. It compounds per cycle only when the loop
is closed.

% ---- References ----------------------------------------------
\begin{thebibliography}{20}
\small

\bibitem{mellhammar2023} Mellhammar, E.; Dahlbom, M.; Vilhelmsson-Timmermand, O.; Strand, S. Tumor Control Probability and Small-Scale Monte Carlo Dosimetry: Effects of Heterogenous Intratumoral Activity Distribution in Radiopharmaceutical Therapy. \textit{Journal of Nuclear Medicine} \textbf{2023}, \textit{64}, 1632-1637. doi:10.2967/jnumed.123.265523.

\bibitem{tamborino2025} Tamborino, G.; Engbers, P.; de Wolf, T. H.; Reuvers, T. G. A.; Verhagen, R.; Konijnenberg, M.; et al. Establishing in vitro dosimetric models and dose--effect relationships for $^{177}$Lu-DOTATATE in neuroendocrine tumors. \textit{Journal of Nuclear Medicine} \textbf{2025}, \textit{66}, 1291--1298. doi:10.2967/jnumed.125.269470.

\bibitem{larouze2023} Larouze, A.; Alcocer-\'Avila, M.; Morgat, C.; Champion, C.; Hindi\'e, E. Membrane and nuclear absorbed doses from $^{177}$Lu and $^{161}$Tb in tumor clusters: effect of cellular heterogeneity and potential benefit of dual targeting --- a Monte Carlo study. \textit{Journal of Nuclear Medicine} \textbf{2023}, \textit{64}, 1619--1624. doi:10.2967/jnumed.123.265509.

\bibitem{perrot2014} Perrot, Y.; Degoul, F.; Auzeloux, P.; Bonnet, M.; Cachin, F.; Chezal, J. M.. Internal dosimetry through GATE simulations of preclinical radiotherapy using a melanin-targeting ligand. \textit{Physics in Medicine \& Biology} \textbf{2014}, \textit{59}, 2183--2198. doi:10.1088/0031-9155/59/9/2183.

\bibitem{tamborino2020} Tamborino, G.; De Saint-Hubert, M.; Struelens, L.; Seoane, D. C.; Ruigrok, E. A. M.; Aerts, A.; et al.. Cellular dosimetry of [$^{177}$Lu]Lu-DOTA-[Tyr$^3$]octreotate radionuclide therapy: the impact of modeling assumptions on the correlation with in vitro cytotoxicity. \textit{EJNMMI Physics} \textbf{2020}, \textit{7}, 12. doi:10.1186/s40658-020-0276-5.

\bibitem{sahota2025} Sahota, S.; Quirindongo, K. A.; Piccolo, J.; Chhabra, A.; Zaid, N.; Deveauxbray, S.; et al.. Enhanced activity localization and microscale dosimetry in alpha-emitter radiopharmaceutical therapy using integrated autoradiography and histological imaging. \textit{Scientific Reports} \textbf{2025}, \textit{15}. doi:10.1038/s41598-025-09277-4.

\bibitem{feijtel2021} Feijtel, D.; Doeswijk, G. N.; Verkaik, N. S.; Haeck, J. C.; Chicco, D.; Angotti, C.; et al.. Inter and intra-tumor somatostatin receptor 2 heterogeneity influences peptide receptor radionuclide therapy response. \textit{Theranostics} \textbf{2021}, \textit{11}. doi:10.7150/thno.51215.

\bibitem{fransson2024} Park, S. W. S.; Fransson, S.; Sundquist, F.; Nilsson, J. N.; Grybäck, P.; Wessman, S.; et al.. Heterogeneous SSTR2 target expression and a novel KIAA1549::BRAF fusion clone in a progressive metastatic lesion following $^{177}$Lutetium-DOTATATE molecular radiotherapy in neuroblastoma: a case report. \textit{Frontiers in Oncology} \textbf{2024}, \textit{14}, 1408729. doi:10.3389/fonc.2024.1408729.

\bibitem{mvogo2026} Mvogo, A.; Essongo, F. E.; Hubert Ben-Bolie, G.. Cancer treatment by radioimmunotherapy: insights from a dynamical model of cancer stem cells and hypoxia effects. \textit{Scientific Reports} \textbf{2026}, \textit{16}, 16975. doi:10.1038/s41598-026-47796-w.

\bibitem{liu2021radcell} Liu, R.; Higley, K. A.; Swat, M. H.; Chaplain, M. A. J.; Powathil, G. G.; Glazier, J. A.. Development of a coupled simulation toolkit for computational radiation biology based on Geant4 and CompuCell3D. \textit{Physics in Medicine \& Biology} \textbf{2021}, \textit{66}, 045026. doi:10.1088/1361-6560/abd4f9.

\bibitem{albers2021} Albers, J.; Svetlove, A.; Alves, J.; Kraupner, A.; di Lillo, F.; Markus, M. A.; et al. Elastic transformation of histological slices allows precise co-registration with microCT data sets for a refined virtual histology approach. \textit{Scientific Reports} \textbf{2021}, \textit{11}. doi:10.1038/s41598-021-89841-w.

\bibitem{ijms2026nanoparticle} Finogenova, Y.; Lipengolts, A.; Klementyeva, O.; Shpakova, K.; Skribitsky, V.; Grigorieva, E.; et al. Linking in vivo imaging to therapeutic outcome with $^{131}$I, $^{177}$Lu, $^{188}$Re nanoparticles. \textit{International Journal of Molecular Sciences} \textbf{2026}, \textit{27}, 7856. doi:10.3390/ijms27177856.

\bibitem{su2023ppgl} Su, D.; Yang, H.; Qiu, C.; Chen, Y.. Peptide receptor radionuclide therapy in advanced Pheochromocytomas and Paragangliomas: a systematic review and meta-analysis. \textit{Frontiers in Oncology} \textbf{2023}, \textit{13}, 1141648. doi:10.3389/fonc.2023.1141648.

\bibitem{liu2025pbt} Liu, R.; H. Swat, M.; A. Glazier, J.; Lei, Y.; Zhou, S.; A. Higley, K.. Developing an Agent-Based Mathematical Model for Simulating Post-Irradiation Cellular Response: A Crucial Component of a Digital Twin Framework for Personalized Radiation Treatment. \textbf{2025}. arXiv:2501.11875.

\bibitem{albanese2026} Albanese, F.; Chiari, G.; Edoardo Delitala, M.. Oxygenation and spatial heterogeneity shape radiotherapy protocol ranking through phenotypic adaptation. \textbf{2026}. arXiv:2606.04004.

\bibitem{cordoni2026} G. Cordoni, F.; Battestini, M.; Missiaggia, M.. A stochastic agent-based extension of the GSM2 model for particle therapy: cell-cycle dynamics, dose-rate dependence, and fractionation effects. \textbf{2026}. arXiv:2604.27630.

\bibitem{zhang2025flash} Zhang, Y.; Huang, C.; Hu, A.; Wang, Y.; Zhu, Y.; Zhou, W.; et al. Unraveling the Redox Mechanisms Underlying FLASH Radiotherapy: Critical Dose Thresholds and NRF2-Driven Tissue Sparing. \textbf{2026}. arXiv:2508.20510.

\bibitem{ryhiner2025tdt} Ryhiner, M.; Song, Y.; Saboury, B.; Glatting, G.; Rahmim, A.; Shi, K.. Mathematical and Computational Nuclear Oncology: Toward Optimized Radiopharmaceutical Therapy via Digital Twins. \textbf{2025}. arXiv:2511.03755.

\bibitem{kinder2026biofilm} Kinder, H.; Faulkner, B.. Modeling radioresistance and radiotropic fitness in multispecies biofilms: a computational framework, implementation audit, and one-way radiation-transport foundation. \textbf{2026}.

\bibitem{li2026chns} Li, C.; Lin, P.; Yu, H.; Zheng, H.. A Thermodynamically Consistent Cahn--Hilliard--Navier--Stokes Model for Tumor Growth. \textbf{2026}. arXiv:2608.06099.

\bibitem{shen2022diffuse} Shen, L.; Lin, P.; Xu, Z.; Xu, S.. Thermodynamically Consistent Diffuse Interface Model for Cell Adhesion and Aggregation. \textbf{2022}. arXiv:2205.07190.

\end{thebibliography}

\end{document}